\documentclass[a4paper,12pt]{article}
\usepackage{amsmath,amssymb}
\usepackage[russian]{babel}
\usepackage[utf8]{inputenc}
\usepackage[T2A]{fontenc}
\usepackage{geometry}
\usepackage{listings}
\usepackage{xcolor}
\usepackage{graphicx}
\usepackage{color}


\geometry{margin=1.5cm}
% Пример темы в стиле тёмной IDE
\definecolor{intj-bg}{HTML}{2B2B2B}
\definecolor{intj-key}{HTML}{CC7832}
\definecolor{intj-str}{HTML}{6A8759}
\definecolor{intj-comm}{HTML}{808080}
\definecolor{intj-id}{HTML}{A9B7C6}
\definecolor{intj-num}{HTML}{6897BB}

\lstdefinestyle{intellij}{
    language=Python,
    backgroundcolor=\color{intj-bg},
    basicstyle=\ttfamily\small\color{intj-id},
    keywordstyle=\color{intj-key}\bfseries,
    stringstyle=\color{intj-str},
    commentstyle=\itshape\color{intj-comm},
    numberstyle=\color{intj-comm}\tiny,
    stepnumber=1,
    frame=single,
    rulecolor=\color{intj-comm},
    showstringspaces=false,
    breaklines=true,
    tabsize=4
}




\begin{document}

% ---------- Титульный лист ----------
    \begin{titlepage}
        \begin{center}
            \textbf{Федеральное государственное автономное образовательное учреждение высшего образования «Национальный исследовательский университет ИТМО} \\
            \vspace{1cm}

            \textbf{Факультет программной инженерии и компьютерной техники} \\
            \vspace{1cm}

            \Large{\textbf{ОТЧЁТ}} \\
            \Large{\textbf{по задачам: A-D}} \\
            \large{по дисциплине «Алгоритмы и структуры данных»} \\
            \vspace{2cm}
        \end{center}

        \begin{flushright}
            \textbf{Выполнил:} \\
            Ануфриев Андрей Сергеевич \\
            P3219 \\
        \end{flushright}

        \vfill
        \begin{center}
            Санкт-Петербург 2026
        \end{center}
    \end{titlepage}

% ---------- Содержание ----------
    \tableofcontents
    \newpage

% ========== ЗАДАНИЕ 1 ==========


    \section{Задание A. Агроном-любитель}

    \subsection*{Листинг кода A}
    \lstinputlisting[style=intellij, label=lst:taskA]{../tasks/A.cpp}

    \subsection*{Суть решения}
    В решении поддерживается скользящее окно, в котором отслеживается текущая длина подряд идущего одинакового элемента (count).\\
    Когда встречаются три одинаковых подряд, левый край сдвигается так, чтобы в окне оставалось не более двух одинаковых подряд.\\
    Для каждого правого конца обновляется длина текущего отрезка и сохраняется лучший найденный отрезок.
    \\
    \textbf{Время: O(n)} — один проход по массиву.\\
    \textbf{Память: O(1)} — используется константное количество доп. памяти.


% ========== ЗАДАНИЕ 2 ==========
    \newpage


    \section{Задание B. Зоопарк Глеба}

    \subsection*{Листинг кода B}
    \lstinputlisting[style=intellij, label=lst:taskB]{../tasks/B.cpp}

    \subsection*{Суть решения}
    Строка длины 2n обрабатывается слева направо; символы двух типов (животные — строчные, ловушки — заглавные) сопоставляются по букве с помощью стека.\\
    В стеке хранятся кортежы (буква в нижнем регистре, индекс в соответствующей группе, is\_animal).\\
    При встрече символа, если верх стека имеет ту же букву и противоположный тип, происходит пара — формируется соответствие животного и ловушки.\\
    В конце проверяется, пуст ли стек (в противном случае — Impossible).
    \\
    \textbf{Время: O(n)} — каждый символ обрабатывается один раз.\\
    \textbf{Память: O(n)} — стек и массив результатов длины O(n).

% ========== ЗАДАНИЕ 3 ==========
    \newpage


    \section{Задание C. Конфигурационный файл}

    \subsection*{Листинг кода C}
    \lstinputlisting[style=intellij,  label=lst:taskC]{../tasks/C.cpp}

    \subsection*{Суть решения}
    Поддерживается хеш-таблица текущих значений переменных (unordered\_map) и стек, в котором для каждого блока\\
    сохраняется вектор пар \{имя, старое\_значение\} для восстановления при выходе из блока.\\
    При присваивании константы — просто сохраняем старое значение в текущем блоке (если находимся внутри блока) и обновляем map.\\
    При присваивании переменной копируем текущее значение rhs и печатаем его (для отладки), \\
    затем аналогично сохраняем и обновляем lhs. При встрече \} восстанавливаем значения из вектора в обратном порядке.
    \\

    \textit{тут сам не справился - запутался с вложенностью и выходом, помогала нейронка}


    \textbf{Время: O(m)} — m строк входа; операции map в среднем O(1).\\
    \textbf{Память: O(v + b)} — где v — число различных переменных, b — максимальная глубина/суммарный объём сохранённых записей (в худшем случае O(m)).


% ========== ЗАДАНИЕ 4 ==========
    \newpage


    \section{Задание D. Профессор Хаос}

    \subsection*{Листинг кода D}
    \lstinputlisting[style=intellij, label=lst:taskD]{../tasks/D.cpp}

    \subsection*{Суть решения}
    Модель процесса симулируется по дням: на каждой итерации текущее число умножается на b, уменьшается на c,\\
    затем ограничивается сверху d. Поскольку k может быть очень большим (до 1e18), \\
    используется обнаружение цикла по состоянию (значению cur): сохраняется отображение value->день и вектор истории.\\
    При повторном появлении состояния вычисляется длина цикла и прямой переход к искомому дню без пошаговой симуляции.\\
    Также обрабатываются случаи завершения эксперимента (cur==0).
    \\

    \textit{пробовал как-то по-умному оптимизировать, но не получилось. В итоге подсказали искать цикл, что он обязательно будет и по нему искать ответ.}


    \textbf{Время: O(min(k, S))} \\
    \textbf{Память: O(min(k, S))}, где S — число возможных различных состояний (ограничено d+1), \\
    т.е. алгоритм работает за линейное по числу уникальных состояний время и память.

\end{document}